% !TeX root = ../tesis.tex

El núcleo criptográfico de \Mithril es un esquema de \textit{stake-based threshold multisignatures} (\STM), cuya meta es convencer a un \textit{verifier} de que una fracción suficiente del stake del sistema firmó un mensaje dado \cite{cryptoeprint:2021/916}. No queremos pedir firmas a todos los participantes activos de la red, sino seleccionar una subcolección de ellos mediante una lotería ponderada por \textit{stake} y luego agregar las firmas individuales ganadoras. En esta sección conviene distinguir dos planos. Para la verificación \STM de certificados, la ruta \Lagrange/Schnorr es conceptualmente central, porque produce la interfaz más natural para circuitos, \textit{public inputs} estructurados y pruebas \ZK. Para la pertenencia de transacciones en \CardanoTransactions, en cambio, el flujo integrado de esta tesis conserva el camino \Pythagoras/\textit{legacy} que hoy expone el \MithrilAggregator productivo.

\subsection{Fases del protocolo}

La documentación oficial organiza \Mithril en tres fases: establecimiento, inicialización y operación \cite{MithrilDocsProtocol}. Consideremos un conjunto fijo de participantes $P = \{1,\ldots,n\}$, donde cada participante $i \in P$ tiene un \textit{stake} $s_i \in \N$. El stake total del sistema es
$$S = \sum_{i=1}^{n} s_i,$$
y el peso relativo del participante $1 \leq i \leq n$ viene dado por $w_i = \frac{s_i}{S}$. El protocolo queda parametrizado por una tripla $(k,m,\phi_f)$, donde $m$ es el número de loterías evaluadas por participante, $k$ es el mínimo de índices ganadores distintos requeridos para la firma agregada y $\phi_f \in (0,1]$ controla la probabilidad de elegibilidad \cite{cryptoeprint:2021/916,MithrilDocsProtocol}.

\paragraph{Fase de establecimiento.}

Se fijan el grupo criptográfico subyacente, los algoritmos de firma, las funciones de hash relevantes y los parámetros $(k,m,\phi_f)$. El resultado de esta fase es una instancia del protocolo
$$\Pi_{\mathsf{Mithril}} = (\mathcal{G}, \Sign, \Verify, k, m, \phi_f),$$
donde $\mathcal{G}$ está compuesto por la familia de grupos y curvas elegidas. En el caso \Lagrange, esta fase fija explícitamente una interfaz Schnorr/\textit{SNARK-friendly}: claves de verificación sobre Jubjub, mensajes de protocolo estructurados y una semántica de elegibilidad diseñada para ser reusada por circuitos.

\paragraph{Fase de inicialización.}

Cada participante $i$ genera material criptográfico propio y registra una entrada asociada a su \textit{stake}. Para la ruta \Lagrange, la información relevante no se agota en una clave pública, sino en una entrada cerrada que ya incorpora la clave Schnorr y el umbral de lotería correspondiente.
$$\ClosedReg_i = \left(\SchnorrVK_i,\LotteryTarget_i,s_i\right)$$
donde $\SchnorrVK_i$ es la clave Schnorr del participante y $\LotteryTarget_i \in \Fq$ es su \texttt{lottery\_target\_value}. Como $\LotteryTarget_i$ depende de $S$, sólo queda bien definido una vez cerrado el registro:
$$\LotteryTarget_i \approx p \left(1-(1-\phi_f)^{w_i}\right),$$

Una vez reunidas todas las entradas cerradas, el registro pasa a un estado autenticado
$$\AuthReg = \{\ClosedReg_1,\ldots,\ClosedReg_n\},$$
del que se deriva una clave de verificación agregada $\AVK = \MerkleRoot{\MerkleTree{R}}$, donde las hojas del árbol de Merkle condensan la información relevante para la ruta SNARK. Esto quiere decir que las hojas de $\MerkleTree{R}$ tienen la forma $\SnarkLeaf{i} = \left(\SchnorrVK_i,\LotteryTarget_i\right)$, y que la \AVK es un \textit{commitment} del comité autenticado para la ronda actual de \STM contra el cual después se verifican las firmas y el Merkle Path asociados.

\paragraph{Fase de operación.}

Dado un mensaje $M$, cada participante $i$ computa una firma individual y evalúa su elegibilidad sobre los índices $\LotteryIndexSet = \{0,\ldots,m-1\}$. La salida local del \textit{signer} puede modelarse como un par
$(\sigma_i, \WonIndexSet_i)$, donde $\WonIndexSet_i \subseteq \LotteryIndexSet$ es el conjunto de índices ganadores para el par $(i, M)$. En el caso \Lagrange, la firma es la \UniqueSchnorrSignature $\sigma_i = (C_i, r_i, c_i)$, pero la implementación no firma directamente una tupla informal compuesta por una raíz de Merkle Tree y el resto del mensaje: primero se deriva $\RigidDigest = \RigidHash(\ProtocolMsg)$, a partir del mensaje \ProtocolMsg usando \ProtocolMessageComputeHash, y recién sobre ese valor se computa la firma. La elegibilidad se determina comparando una evaluación derivada de esa firma contra el \texttt{lottery\_target\_value} $\LotteryTarget_i$ ya comprometido en el registro cerrado.

La condición de elegibilidad para cada índice $j \in \LotteryIndexSet$ puede escribirse como un predicado $\mathsf{Win}(i,j,M) \in \{0,1\}$, y entonces $\WonIndexSet_i = \{j \in \LotteryIndexSet: \mathsf{Win}(i, j, M) = 1 \}$. El \textit{aggregator} recibe los pares $(\sigma_i, \WonIndexSet_i)$, descarta aquellos que no verifican, deduplica índices repetidos y únicamente puede producir una firma agregada válida si el conjunto total de índices ganadores distintos satisface

$$\left| \bigcup_{i=1}^{n} \WonIndexSet_i \right| \geq k$$.

Observamos entonces que no alcanza con acumular muchas firmas si en realidad están justificando los mismos índices. El protocolo \STM produce una firma agregada que certifica que un subconjunto autenticado del comité, firmó el mensaje $M$ con al menos $k$ índices válidos y distintos.

\subsection{Registro cerrado y clave de verificacion agregada SNARK}

Una vez fijado el stake total $S$, cada participante queda representado por una entrada cerrada que incorpora, ademas del \textit{stake}, una clave Schnorr y el umbral de lotería correspondiente. Si denotamos esa entrada por $\ClosedReg_i$, la parte relevante para la \AVK es $\SnarkLeaf{i} = \bigl(\SchnorrVK_i,\LotteryTarget_i\bigr)$,
donde $\SchnorrVK_i$ es la clave de verificación Schnorr del participante y $\LotteryTarget_i$ es el \texttt{lottery\_target\_value} derivado de $(\phi_f,s_i,S)$. El código calcula $\LotteryTarget_i$ durante el cierre del registro como una aproximación (en el cuerpo base de Jubjub $\Fq$), de
$$p \cdot \left(1-(1-\phi_f)^{w_i}\right),$$
donde $w_i = \frac{s_i}{S}$. La serialización de la hoja es
$$\mathsf{enc}(\SnarkLeaf{i}) = \mathsf{enc}_{64}(\SchnorrVK_i) \concat \mathsf{enc}_{32}(\LotteryTarget_i),$$
de donde se obtiene la hoja $\ell_i = \SnarkLeafDigest(\mathsf{enc}(\SnarkLeaf{i}))$. Entonces, la \AVK está dada por la raíz del Merkle Tree del conjunto de hojas así como el \textit{stake} total:
$$\mathsf{AVK} = \left(\SnarkRoot, S \right), \qquad \mathcal{L} = \{\ell_i\}_{i=1}^{n},$$
donde las hojas se ordenan canónicamente antes de armar el árbol, primero por $\LotteryTarget_i$ y luego por $\SchnorrVK_i$ \cite{MithrilSTMRepo}.

\subsection{Esquema de firmas en \Lagrange}

En esta era de \Mithril, se produce para cada participante una clave de firmado $\SchnorrSK_i$ y su correspondiente clave de verificación $\SchnorrVK_i$. Este es el camino más importante para la verificación \STM de certificados del \zkBridge, ya que para obtener argumentos sucintos verificables \Onchain en \ChainB no alcanza con transportar la \BLS \textit{multisignature} de la era \textit{legacy}/\Pythagoras. Esto no significa que toda la evidencia del bridge haya migrado a \Lagrange: la pertenencia de \CardanoTransactions en el prototipo integrado se trata por separado mediante el camino \Pythagoras actualmente disponible.

Dentro de este camino, aparecen dos variantes de firma Schnorr. La primera es la \StandardSchnorrSignature, formada solo por un \textit{challenge} $c$ y su correspondiente respuesta $r$. Según la propia implementación, esta variante se usa en el \SNARK \GenesisCertificate. Su desafio se recomputa como un hash \PoseidonHash del \textit{domain separation tag}, la clave de verificación, el punto aleatorio y el mensaje. La segunda es la \UniqueSchnorrSignature, que agrega un \CommitmentPoint determinista dependiente solo del mensaje y de la clave Schnorr \SchnorrSK. Ese punto permite evitar \textit{grinding attacks} sobre la lotería de \MithrilAggregator, y por lo tanto es la forma que el circuito \HaloTwo reutiliza después \cite{MithrilSTMRepo}.

\begin{algorithm}[ht]
	\caption{\StandardSchnorrSignature en la era \Lagrange}
	\label{alg:mithril-lagrange-standard-schnorr-sign}
	\begin{algorithmic}[1]
		\Require Mensaje $\mathbf{M} = (m_1, \ldots, m_t)$, clave secreta Schnorr $\sk^{\mathsf{sch}}$
		\Ensure $\sigma = (r, c)$
		\State $\mathsf{vk} \gets \sk_{\mathsf{sch}} \cdot G$
		\State $\rho \gets \texttt{RandomNonZeroScalar}()$
		\State $R \gets \rho \cdot G$
		\State $c \gets \texttt{Poseidon}(\texttt{DST}_{\mathsf{std}} \concat \mathsf{vk} \concat R \concat \mathbf{m})$
		\State $r \gets \rho - c \cdot \sk_{\mathsf{sch}}$
		\State \Return $(r, c)$
	\end{algorithmic}
\end{algorithm}

\begin{algorithm}[ht]
	\caption{\UniqueSchnorrSignature en la era \Lagrange}
	\label{alg:mithril-lagrange-unique-schnorr-sign}
	\begin{algorithmic}[1]
		\Require Mensaje $\mathbf{M} = (m_1,\ldots,m_t)$, clave Schnorr $\sk^{\mathsf{sch}}$
		\Ensure $\sigma = (C, r, c)$
		\State $\mathsf{vk} \gets \sk_{\mathsf{sch}} \cdot G$
		\State $H \gets \texttt{HashToCurve}(\mathbf{m})$
		\State $C \gets \sk_{\mathsf{sch}} \cdot H$ \Comment{\CommitmentPoint}
		\State $\rho \gets \texttt{RandomNonZeroScalar}()$
		\State $R_1 \gets \rho \cdot H$
		\State $R_2 \gets \rho \cdot G$
		\State $c \gets \texttt{Poseidon}(\texttt{DST}_{\mathsf{uniq}} \concat H \concat \mathsf{vk} \concat C \concat R_1 \concat R_2)$
		\State $r \gets \rho - c \cdot \sk^{\mathsf{sch}}$
		\State \Return $(C, r, c)$
	\end{algorithmic}
\end{algorithm}

\begin{algorithm}[ht]
	\caption{Firma de un participante con chequeo de lotería en \Lagrange}
	\label{alg:mithril-lagrange-snark-signer-sign}
	\begin{algorithmic}[1]
		\Require Mensaje $\ProtocolMsg$, clave secreta Schnorr $\SchnorrSK_i$, \texttt{lottery\_target\_value} $\LotteryTarget_i$, parámetro $m$
		\Ensure Par $(\sigma_i, \WonIndexSet_i)$ con $\sigma_i = (C_i,r_i,c_i)$ y $\WonIndexSet_i \subseteq \LotteryIndexSet=\{0,\ldots,m-1\}$, o \LotteryLost
		\State $\RigidDigest \gets \RigidHash(\ProtocolMsg)$ \Comment{\ProtocolMessageComputeHash}
		\State $\sigma_i = (C_i, r_i, c_i) \gets \UniqueSchnorrSignature(\RigidDigest,\SchnorrSK_i)$
		\State $\WonIndexSet_i \gets \{\, j \in \LotteryIndexSet : \LotteryEval(\sigma_i, j) \leq \LotteryTarget_i \,\}$
		\If{$\WonIndexSet_i = \emptyset$}
			\State \Return \LotteryLost
		\Else
			\State \Return $(\sigma_i, \WonIndexSet_i)$
		\EndIf
	\end{algorithmic}
\end{algorithm}

El Algoritmo \ref{alg:mithril-lagrange-snark-signer-sign} está formulado para reflejar la semántica observable en el código de \Mithril. En particular, el participante no firma una Merkle Root aislada ni devuelve una firma ``a medio terminar'' para que el \textit{aggregator} elija luego los índices. La secuencia correcta es: derivar primero el \textit{digest} rígido \RigidDigest del mensaje, producir después una \UniqueSchnorrSignature sobre ese valor y, finalmente, calcular en forma local el conjunto $\WonIndexSet_i$ de índices de loterías ganadas. La notación \(\LotteryEval(\sigma_i, j)\) se refiere a la derivación determinista previamente expuesta.

\subsection{Firmas individuales y agregación}

Durante la fase de operación, cada participante evalúa para un mensaje \ProtocolMsg y para cada índice $j \in \LotteryIndexSet=\{0,\ldots,m-1\}$ si gana la lotería correspondiente a ese índice. Si resulta elegible para uno o mas indices, genera y envía firmas individuales asociadas a esos índices. La agregacion posterior no presupone honestidad del \textit{aggregator}: cualquier tercero con acceso a las firmas individuales puede verificar que son correctas y filtrar aquellas que no satisfacen los chequeos de elegibilida, así como hacer la agregación \cite{MithrilDocsProtocol}. Esta propiedad es útil para el \zkBridge, porque evita introducir un supuesto adicional de confianza sobre el \MithrilAggregator.

\paragraph{Agregacion en \Lagrange.}
En la ruta SNARK, el \MithrilAggregator recibe un conjunto de firmas
$$\SnarkSignatureSet = \{(\sigma_i^{\mathsf{sch}}, \WonIndexSet_i, i)\}_{i \in P'},$$
del conjunto de participantes $P' \subseteq P$ que le enviaron firmas sobre el mensaje hasheado $\RigidDigest = \RigidHash(\ProtocolMsg).$

El primer paso consiste en descartar las firmas que no verifiquen respecto de la \SchnorrVK registrada para cada participante, o cuyos índices reclamados no coincidan con los que se recomputan a partir de $\LotteryTarget_i$. Este cómputo usa la información de la hoja $\SnarkLeaf{i} = (\SchnorrVK_i,\LotteryTarget_i)$, que luego será enlazada al Merkle tree comprometido por \(\SnarkAVK\). El resultado es un conjunto de candidatos válidos $\SnarkValidSignatureSet \subseteq \SnarkSignatureSet$.

Recordemos que en \STM el umbral se mide en índices de lotería distintos, no en cantidad de firmas, para poder cubrir al menos $k$ valores distintos de $\LotteryIndexSet$. Si definimos
$$\CandidateIndexSet = \bigcup_{(\sigma_i^{\mathsf{sch}},\WonIndexSet_i,i)\in \SnarkValidSignatureSet} \WonIndexSet_i,$$
la agregación solo puede avanzar cuando \(|\CandidateIndexSet| \ge k\). Una vez superado ese umbral, el código debe elegir qué índices utilizar, pero sin beneficiar o perjudicar a ningún participante específico. Primero, deriva una \textit{seed} pública $\SnarkSeed = \SHATwoFiveSix(\texttt{DST}_{\mathsf{seed}} \concat \RigidDigest)$. Para seleccionar que indices sobreviven, calcula para cada índice candidato $j$ un puntaje
$$h_{\mathsf{idx}}(j) =
\SHATwoFiveSix(\texttt{DST}_{\mathsf{idx}} \concat \SnarkSeed \concat j),$$
y conserva exactamente los $k$ indices con menor valor. Denotamos por $\SelectedIndexSet $ a ese conjunto, con
$|\SelectedIndexSet| = k$ y $\SelectedIndexSet \subseteq \CandidateIndexSet$.

Como para cada $j \in \SelectedIndexSet$ puede haber varios participantes que lo hayan ganado, el \MithrilAggregator hace un desempate, que depende de una segunda familia de hashes:
$$h_{\mathsf{dedup}}(j,i) =
\SHATwoFiveSix(\texttt{DST}_{\mathsf{dedup}} \concat \SnarkSeed \concat j \concat i),$$
de modo que el participante retenido para el indice $j$ es
$$i^{\star}(j) = \arg\min\{h_{\mathsf{dedup}}(j,i) : j \in \WonIndexSet_i\}.$$

Esta selección deja un mapa ordenado $\SnarkSelectionMap : \SelectedIndexSet \to P'$, con exactamente un participante por índice y exactamente $k$ indices retenidos.

El paso final produce un \textit{witness} $w \in \WitnessSpace$ para el circuito de \HaloTwo de \Mithril. Para cada $j \in \SelectedIndexSet$, el \MithrilAggregator construye
$$\SnarkWitness_j = \left(\SnarkLeaf{i^{\star}(j)}, \mathsf{Path}_{i^{\star}(j)}, \sigma^{\mathsf{sch}}_{i^{\star}(j)}, j\right),$$

donde $\mathsf{Path}_{i^{\star}(j)}$ es el Merkle Path de $\SnarkLeaf{i^{\star}(j)}$ hasta $\SnarkRoot$. Entonces,
$$\SnarkWitness =
\left\{\left(\SnarkLeaf{i^{\star}(j)}, \mathsf{Path}_{i^{\star}(j)}, \sigma^{\mathsf{sch}}_{i^{\star}(j)}, j\right)\right\}_{j \in \SelectedIndexSet} \in \WitnessSpace.$$

Posteriormente, el argumento construido con \HaloTwo va a demostrar que todas esas entradas son coherentes respecto de la $\SnarkAVK$, del mismo hash $\RigidDigest$ y de la condición de quórum $|\SelectedIndexSet| = k$. Es por esto que esta subsección se enfoca en mostrar cómo todo puede ser expresado de forma determinística, para llevar a un circuito con menos dificultades.
